% Printed212; principalPDF219.
\par\clearpage
\noindent212\hfill\textsc{Briefwechsel.}\hfill\mbox{}\par\medskip
\begin{center}Konzise Musterrechnung.\\Neuwerk.\end{center}
\[
\begin{array}{r@{\qquad}r@{\qquad}r}
-266575{,}038&+95076{,}254&\\[8pt]
\begin{array}[t]{r}
133{,}287\,519\\10{,}663\,002\\\hline
143{,}950\,521\\2^\circ\,23'\,57''{,}03\\51\quad23\quad20{,}61\\\hline
53\quad47\quad17{,}64\\8\quad15{,}73\\\hline
53\quad55\quad33{,}37\\-31{,}40\\\hline
53^\circ\,55'\quad1''{,}97\\\hline
\end{array}&
\begin{array}[t]{r}
4{,}978\,0721\\8{,}508\,8620\\\hline
3{,}486\,9341\\9{,}769\,9904\\0{,}137\,5588\\9{,}956\,1442\\1{,}403\,2842\\\hline
3{,}716\,9437\\-764\\\hline
1{,}496\,9872\\-850\\\hline
3{,}624\,4929\\-923
\end{array}&
\begin{array}[t]{r}
7{,}43389\\\hline5{,}44934\\5{,}49566\\5{,}53146\\[8pt]
5210''{,}355\\1^\circ\,26'\,50''{,}355\\\hline
4211''{,}149\\1^\circ\,10'\,11''{,}15\\\hline
\end{array}
\end{array}
\]
\bigskip
\begin{center}\textsc{Gauss} an \textsc{Schumacher}.\quad Göttingen, 30. April 1830.\end{center}
. . . . . . Ich fahre jetzt mit dem fort, was die Hilfstafeln betrifft, heute nur mit den Formeln für die übrigen drei Spalten, die bei der Berechnung der Länge, Breite und Konvergenz gebraucht werden. Ich weiß aber nicht, ob ich sie in meinem vorigen Brief mit \(D,E\) und \(F\) bezeichnet habe.

Zahlen, deren Logarithmen angesetzt werden:
\[
\begin{array}{@{}l@{\;}l@{}}
\text{Länge . . . . }[D=]&h\left(\frac34-\frac14\cos2\varphi+\frac{ee}{2(1-ee)}\cos\varphi^4\right)\\[6pt]
\text{Breite . . . . }[E=]&h\left(\frac34+\frac{2-11ee}{4(1-ee)}\cos\varphi^2+\frac{5ee}{2(1-ee)}\cos\varphi^4-\frac{e^4}{(1-ee)^2}\cos\varphi^6\right)\\[6pt]
\text{Konvergenz }[F=]&h\left(1-\frac{ee}{1-ee}\cos\varphi^4-\frac{2e^4}{(1-ee)^2}\cos\varphi^6\right).
\end{array}
\]
\[
\text{Hier ist }h=\frac{k\,.\,10^7}{3\,.\,(206265)^2};\qquad
\begin{array}{l}
k=\text{Modul}=0{,}434\rlap{\raisebox{-.5ex}{\tiny.}}^{\prime}2945\\
\log h=5{,}531\,8128-10.
\end{array}
\]

% Printed213; principalPDF220.
\par\clearpage
\noindent\mbox{}\hfill{\small\textsc{Über die Formeln für die Hannoversche Landesvermessung.}}\hfill213\par\medskip
Es ist zu bemerken, dass diese Formeln vollständig sind, das heißt, sie sind keine unendlichen Reihen, sondern bestehen nur aus diesen Gliedern. . . . . . .

Um Ihr Vertrauen zu \textsc{Schmidts} Rechnung zu vergrößern, bemerke ich, dass er die zwei Hauptelemente der Erddimensionen viermal berechnet hat, aber nur einmal hat er wegen eines Rechenfehlers von Neuem gerechnet. Nämlich
\begin{list}{}{\setlength{\leftmargin}{2em}\setlength{\itemsep}{0pt}\setlength{\parsep}{0pt}}
\item[1)] Zahlen in meiner Breitenbestimmung etc.

Diese hatten einen Rechenfehler enthalten, den er später verbesserte, daher
\item[2)] die Zahlen in seiner Geographie und in Ihren A. N.

Erst später machte ich ihn auf die Korrektur der in Ostindien gebrauchten Maßstäbe aufmerksam, daher die
\item[3)] Rechnung, deren Resultat in der Vorrede des Buches steht.

Endlich hat er seitdem eine vierte Rechnung gemacht, nicht wegen eines Rechenfehlers, sondern um die ihm erst nachher bekannt gewordenen Resultate von \textsc{Struves} Gradmessung mit unter die Daten aufzunehmen. Das Resultat
\item[4)] ist mir von ihm handschriftlich mitgeteilt und dasselbe, das meinen neuen Hilfstafeln zugrunde liegt, nämlich
\end{list}
\[
\text{Abplattung}=\frac1{297{,}732};\qquad
\frac{\text{Erdquadrant}}{90}=57008{,}551.
\]
. . . . . .
\bigskip
\begin{center}\textsc{Gauss} an \textsc{Schumacher}.\quad Göttingen, 17. Mai 1831.\end{center}
. . . . . . Es gibt \textls[120]{mehrere} Wege, um aus den Breiten und Längen die Koordinaten zu berechnen, die jeder an seiner Stelle ihre eigentümlichen Vorzüge haben. Für den Fall, für den Sie solche Rechnungen benötigen, ist es am bequemsten, sich der Reihen zu bedienen, die dann so schnell konvergieren, dass sehr wenige Glieder ausreichen.

Entsprechen die Koordinaten \(x,y\) der Breite \(\varphi\) und Länge \(\lambda\), so ist die Form diese\textsuperscript{[*)]}:
\[
\begin{array}{l}
x=A-A'\lambda\lambda-A''\lambda^4-\text{etc.}\\
y=B\lambda+B'\lambda^3+\text{etc.},
\end{array}
\]
\par\noindent\rule{4cm}{.4pt}\par
{\footnotesize\noindent\hspace*{\parindent}[*) Vergl. Art. 13, S. 158.]\par}

% Printed214; principalPDF221.
\par\clearpage
\noindent214\hfill\textsc{Briefwechsel.}\hfill\mbox{}\par\medskip
\noindent wo \(A,A',A''\), etc. und \(B,B'\), etc. Funktionen von \(\varphi\) sind, deren numerische Werte man für diejenigen runden Grade (oder halben Grade etc.) zu berechnen hat, die innerhalb der Karte vorkommen.

Zur Berechnung von \(A\) bedürfen Sie meiner Anleitung nicht, da sie lediglich von der Rektifikation eines elliptischen Bogens abhängt. Ist nämlich \(\varphi^\circ\) die Breite desjenigen Ortes, von dem aus man die Koordinaten \(x\) nach Süden zählt, \(u^\circ\) dessen wirkliche Entfernung vom Äquator und \(u\) die variable Äquatorentfernung des Parallelkreises \(\varphi\) (beide, \(u,u^\circ\), in derjenigen Einheit ausgedrückt, die man für die Koordinaten gewählt hat), so ist \(A=u^\circ-u\). Sie können dazu auch eine Ihnen schon früher mitgeteilte Hilfstafel benutzen\textsuperscript{[*)]}. Die Werte der anderen Koeffizienten können Sie nach folgenden Formeln berechnen, die \textls[120]{absolut} genau sind.
\[
\begin{aligned}
A'&=\frac{asc}{2p}\\[4pt]
A''&=\frac{asc}{24(1-ee)^2p}\bigl(5-ee-(6+6ee)ss+(9ee+3e^4)s^4-4e^4s^6\bigr)\\[4pt]
B&=\frac{ac}{p}\\[4pt]
B'&=\frac{ac}{6(1-ee)p}(1-2ss+ees^4).
\end{aligned}
\]
Wobei Folgendes zu bemerken ist.

Es bedeuten:
\begin{quote}
\(a\) den Halbmesser des Äquators, \(e\) die Exzentrizität;
\[
\GaussRightBrace{\begin{array}{l}
c=\cos\varphi\\s=\sin\varphi\\p=\surd(1-eess)
\end{array}}\ \text{Kürze halber;}
\]
\end{quote}
\noindent ferner ist für \(\lambda\) der Bogen \(57^\circ17'45''\) als Einheit angenommen. Will man also etwa zuletzt die Koordinaten von Grad zu Grad für \(\lambda\) berechnen, so empfiehlt es sich, gleich anfangs dem Koeffizienten \(B\) den Faktor \(\frac1{57{,}296}=\frac{3600}{206265}=\frac\pi{180}\) beizufügen und ebenso die Koeffizienten \(A',B',A''\) sogleich mit der zweiten, dritten, vierten Potenz von \(\frac\pi{180}\) zu multiplizieren.

Die Formeln fallen etwas einfacher aus, wenn man \(\frac{ee}{1-ee}=m\) setzt und sie danach umformt. In dieser Form habe ich sie selbst zu meinen
\par\noindent\rule{4cm}{.4pt}\par
{\footnotesize\noindent\hspace*{\parindent}[*) Seite 206.]\par}

% Printed215; principalPDF222.
\par\clearpage
\noindent\mbox{}\hfill{\small\textsc{Über die Formeln für die Hannoversche Landesvermessung.}}\hfill215\par\medskip
\noindent Rechnungen angewandt, aber nicht aufgehoben\textsuperscript{[*)]}. Ich überlasse also die sehr kleine Arbeit jener Umformung Ihnen selbst. Indem ich \(\varphi^\circ=51^\circ31'47''{,}85\) annehme, für die Abplattung Herrn \textsc{Schmidts} letzte Bestimmung zugrunde lege und als Längeneinheit den \(10\,000\,000.\) Teil des Erdquadranten in dieser Gestalt wähle, finde ich, \(\lambda=n\) Grad gesetzt:
\[
\begin{array}{c@{\quad}l@{\qquad}l}
\varphi&x=&y=\\
51^\circ&+58947{,}1-475{,}95nn-0{,}0167n^4&70180{,}0n-0{,}737n^3\\
52^\circ&-52287{,}0-472{,}16nn-0{,}0154n^4&68660{,}7n-0{,}840n^3\\
53^\circ&-163539{,}8-467{,}79nn-0{,}0140n^4&67120{,}2n-0{,}936n^3\\
54^\circ&-274811{,}2-462{,}85nn-0{,}0127n^4&65559{,}1n-1{,}026n^3\\
55^\circ&-386100{,}9-457{,}34nn-0{,}0114n^4&63977{,}8n-1{,}109n^3.
\end{array}
\]
. . . . . .
\bigskip
\begin{center}\textsc{Gauss} an \textsc{Schumacher}.\quad Göttingen, 25. Juni 1831.\end{center}
. . . . . . Was die Berechnung der Koordinaten betrifft, so kommt es auf zwei Aufgaben an, nämlich:

1) Aus der wirklichen Länge \(R\) einer kürzesten Linie auf dem Sphäroid, deren Endpunkte in der Darstellung jeweils die Koordinaten \(x,y;\ x',y'\) haben, die Entfernung [in] der Darstellung, d. i. die Größe \(r=\surd((x'-x)^2+(y'-y)^2)\) zu finden. Hier ist die Lösung der umgekehrten Aufgabe. Für alle Ihre Fälle gilt mit hinreichender Genauigkeit
\[
\log R=\log r-k\,.\,\frac{(1-ee\sin\varphi^2)^2}{2aa(1-ee)}\,.\,\frac{yy+yy'+y'y'}3.
\]
Da man dabei \(y\) und \(y'\) nicht sehr genau zu kennen braucht, reicht die vorläufige Berechnung der Koordinaten aus, die man so vornimmt, als ob alles in plano wäre. So dient die Formel auch für die Aufgabe, \(r\) aus \(R\) zu finden. \(k\) ist der Modul der Briggs'schen Logarithmen.

Zur wirklichen Berechnung setzte ich die Formel in folgende Gestalt:
\[
\log r=\log R+\{\alpha(y+y')^2+\beta(y-y')^2\}q,
\]
\noindent wo
\[
\alpha=\frac{k}{4\,.\,206265},\quad\beta=\frac{k}{12\,.\,206265},\quad q=206265\,.\,\frac{(1-ee\sin\varphi^2)^2}{2aa(1-ee)};
\]
\par\noindent\rule{4cm}{.4pt}\par
{\footnotesize\noindent\hspace*{\parindent}[*) Diese Formeln findet man in Art. 13, S. 157 unten.]\par}

% Printed216; principalPDF223.
\par\clearpage
\noindent216\hfill\textsc{Briefwechsel.}\hfill\mbox{}\par\medskip
\noindent hier ist
\[
\begin{array}{l}
\log\alpha=3{,}72130-10\\
\log\beta=3{,}24418-10.
\end{array}
\]
Für \(q\) habe ich eine Hilfstafel, die bloß von \(x\) abzuhängen braucht und die ich Ihnen, wie ich glaube, schon einmal mitgeteilt habe\textsuperscript{[*)]}. Offenbar sind \(\alpha,\beta\) von der Maßeinheit unabhängig, aber \(q\) ist es nicht. Eine Hilfstafel dafür hängt sowohl von der Maßeinheit als auch vom Anfang der \(x\) ab. Meine müssten Sie also, wenn Sie eine andere Einheit und einen anderen Anfangspunkt brauchen, erst transformieren.

Übrigens erhalten Sie so \(\log r-\log R\) als Dezimalbruch. Wollen Sie den Wert gleich in Einheiten der 7. Dezimalstelle haben, so brauchen Sie nur \(\log\alpha=0{,}72130,\ \log\beta=0{,}24418\) anzuwenden.

2) Aus einem im Punkt \(P\) gemessenen Winkel zwischen \(PP'\) und \(PP''\) den Winkel zu finden, der ihm in der Darstellung entspricht. Auch hier tritt wieder die umgekehrte Aufgabe an die Stelle.
\begin{center}
\begin{tikzpicture}[x=1cm,y=1cm,font=\normalsize]
\node at (1.8,2) {Sphäroid};\node at (9,2) {Zeichnung in plano};
\draw (0,.7) circle(.035);\node[above left] at (0,.7) {$P''$};
\draw (2.1,1.2) circle(.035);\node[above right] at (2.1,1.2) {$P$};
\draw (2.85,.25) circle(.035);\node[above right] at (2.85,.25) {$P'$};
\draw (7.1,.7) circle(.035);\node[above left] at (7.1,.7) {$p''$};
\draw (9.3,1.2) circle(.035);\node[above right] at (9.3,1.2) {$p$};
\draw (10.05,.25) circle(.035);\node[above right] at (10.05,.25) {$p'$};
\end{tikzpicture}
\end{center}
Es seien die Koordinaten von \(p,p',p''\) der Reihe nach
\[
\begin{array}{ccc}x,&x',&x'',\\y,&y',&y'',\end{array}
\]
\noindent dann ist:
\[
\begin{aligned}
P'PP''={}&p'pp''-q^*(x'-x)\left(\frac{2y+y'}3\right)\\
&+q^{**}(x''-x)\left(\frac{2y+y''}3\right).
\end{aligned}
\quad\text{(gibt das Gesuchte in Bogensekunden)}
\]
Sie können für \(q^*,q^{**}\) den Wert von \(q\) anwenden, der dem Argument \(x\) entspricht. Wollen Sie genauer vorgehen, so ist für \(q^*\) das Argument \(\frac{2x+x'}3=x+\frac13(x'-x)\) und für \(q^{**}\) das Argument \(\frac{2x+x''}3=x+\frac13(x''-x)\). Ich selbst setze immer \(\frac{2y+y'}3\) in die Form \(y+\frac13(y'-y)\), etc. Sie sehen, dass dann das ganze Verfahren darauf hinausläuft, erst jedem gemessenen Azimut (von
\par\noindent\rule{4cm}{.4pt}\par
{\footnotesize\noindent\hspace*{\parindent}[*) Seite 210, wo \(H=q\) ist.]\par}

% Printed217; principalPDF224.
\par\clearpage
\noindent\mbox{}\hfill{\small\textsc{Über die Formeln für die Hannoversche Landesvermessung.}}\hfill217\par\medskip
\noindent \(P\) nach \(P'\)) die Korrektur \(-q(x'-x)(y+\frac13(y'-y))\) beizufügen, woraus das entsteht, was ich Azimut in plano nenne. Unter Azimut auf dem Sphäroid verstehe ich hier aber nicht das astronomische Azimut, d. i. nicht den Winkel mit dem wirklichen Meridian, sondern mit einer Linie, die dem Fundamentalmeridian parallel ist, oder strenger, mit einer Linie auf dem Sphäroid, die in der Darstellung in plano eine Parallele zur Abszissenlinie ergibt.

Azimut in plano ist also immer \(\operatorname{arc\,tang}\frac{y'-y}{x'-x}\).

Nachdem jene Korrekturen angewandt sind, hat man mit allen Winkeln so zu rechnen, als ob alles in plano wäre, und wie dann die Koordinaten zu berechnen sind, darüber bedürfen Sie keiner Vorschriften. Offenbar ist auch für diese Rechnung eine näherungsweise Kenntnis von \(x,y,x'\), etc. hinreichend, wie man sie erhält, wenn man anfangs ohne Korrektur rechnet. Will man alles bis auf \(0''{,}001\) in Übereinstimmung haben, so kann man die Rechnung nötigenfalls nachbessern, nachdem \(x,y\), etc. genauer bekannt sind.

Übrigens sehen Sie leicht, dass die obige Rechnung 1), d. i. der Übergang von \(R\) auf \(r\), nur bei einer Linie vorzunehmen ist (der Basis). Danach ist die Kenntnis der einzelnen Linien auf dem Sphäroid allgemein gesprochen ohne Interesse, es sei denn, dass man wieder zu einer neuen Basis gelangt, wo der umgekehrte Weg (von \(r\) nach \(R\)) anzuwenden ist. Ohnehin sind in allen Ihren Fällen \(R\) und \(r\) immer sehr nahe beieinander, jedenfalls hat der Unterschied auf das Zentrieren der Winkel keinen merklichen Einfluss. Im westlichsten Teil von Westfalen habe ich zwar darauf Rücksicht genommen, aber bloß zur Ehre der Rechnung, denn tatsächlich bringt es auch da nichts.
\begin{center}\textsc{Gauss} an \textsc{Schumacher}.\quad Göttingen, 9. Dezember 1838.\end{center}
Die verlangte Formel ist folgende\textsuperscript{[*)]}:
\[
\begin{aligned}
\varphi={}&\psi+(6ff+48f^4+426f^6+4080f^8\ldots)\sin2\psi\\
&\quad+(21f^4+336f^6+4264f^8\ldots)\sin4\psi\\
&\qquad+\left(\frac{302}3f^6+2416f^8\ldots\right)\sin6\psi\\
&\qquad\quad+\left(\frac{1097}2f^8\ldots\right)\sin8\psi\\
&\qquad\qquad\text{etc.}
\end{aligned}
\]
\par\noindent\rule{4cm}{.4pt}\par
{\footnotesize\noindent\hspace*{\parindent}[*) Vergl. Art. 25, S. 179.]\par}
\noindent{\small IX.}\hfill28\hspace*{.12\linewidth}\par

% Printed218; principalPDF225.
\par\clearpage
\noindent218\hfill{\footnotesize\textsc{Briefwechsel.\quad Über die Formeln für die Hannoversche Landesvermessung.}}\par\medskip
Da ich mich seit ziemlich vielen Jahren nicht mit diesen Dingen beschäftigt habe, hat es mich erst viel Suchen gekostet, bis ich die Formel wiedergefunden habe, und nachher erneutes Suchen, oder vielmehr Rechnen, um die Bedeutung des \(f\) zu ermitteln, die auf dem Blatt nicht vermerkt war. Es ist aber \(f=\frac14e\), wobei \(e\) die Exzentrizität der erzeugenden Ellipse bedeutet. Oder wenn die Abplattung \(\frac1n\) ist, d. i. \(\frac{\text{Erdachse}}{\text{Durchm. d. Äquators}}=\frac{n-1}n\), so ist \(f=\frac{\surd(2n-1)}{4n}\). Ich habe dieses \(f\) deswegen gebraucht, weil, wenn ich \(e\) gebraucht hätte, sämtliche Koeffizienten Brüche geworden wären, da sie hier größtenteils ganze Zahlen sind. Die numerische Rechnung hatte ich durchgehend mit Logarithmen auf 10 Dezimalstellen geführt.
\begin{center}\rule{3cm}{.4pt}\end{center}
\bigskip
\begin{center}\textsc{Bemerkungen.}\end{center}
\begingroup\footnotesize
Die ersten beiden und der letzte der vorstehenden Briefe sind nach den im \textsc{Gauss}-Archiv befindlichen Originalen abgedruckt, während für die Briefe vom 17. Mai und 25. Juni 1831 nach den Originalen angefertigte Kopien benutzt wurden. Im Abdruck sind einige Schreib- und kleinere Rechenfehler berichtigt worden. Der Brief vom 9. Dezember 1838 ist die Antwort auf eine Bitte \textsc{Schumachers}, ihm die Ableitung der im Brief vom 18. April 1830 gegebenen numerischen Formel für \(\varphi-\psi\) mitzuteilen. Die \textsc{Gauss}sche Angabe
\[
\varphi=\psi+520''{,}469\,3364\sin2\psi\ \text{etc.}
\]
\noindent enthält jedoch einen Rechenfehler, wie auch das noch vorhandene Blatt zeigt, das \textsc{Gauss} zur Rechnung benutzt hatte. Es muss heißen, wie S. 206 angegeben ist:
\[
\varphi=\psi+520''{,}463\,3364\sin2\psi\ \text{etc.}
\]
\noindent Dadurch wird auch die letzte Stelle in der dort gegebenen Tabelle für \(\varphi-\psi\) unsicher. Ferner erhält man für Göttingen \(\psi^\circ=51^\circ23'20''{,}6082\), S. 207, anstelle von \(\ldots20''{,}6024\), wie es im Original heißt.
\par\noindent\hfill\textsc{Krüger, Börsch.}\par
\endgroup
\vspace{1.8cm}\begin{center}\rule{4cm}{.4pt}\end{center}

% Unnumbered title219; principalPDF226.
\par\clearpage\thispagestyle{empty}
\null\vspace*{.23\textheight}
\begin{center}{\Large TRIGONOMETRISCHE PUNKTBESTIMMUNG.}\end{center}
\vfill\noindent\hfill28\textsuperscript{*}\hspace*{.15\linewidth}\par\vspace*{.14\textheight}

% Blank220; principalPDF227.
\par\clearpage\null\thispagestyle{empty}

% Unnumbered opening221; principalPDF228.
\par\clearpage\thispagestyle{empty}
\vspace*{2.8\baselineskip}
\begin{center}{\Large\textls[160]{NACHLASS.}}\end{center}
\vspace{2\baselineskip}
\begin{center}[1.]\\[5pt]Endergebnis für den Ort eines Punktes in einer Ebene, der von drei bekannten\\aus angeschnitten ist.\end{center}
Es bedeuten \(10,20,30\) die drei beobachteten Richtungen [nach \(P\)] und \(\alpha,\beta,\gamma\) die entsprechenden Entfernungen.

Die drei einzelnen Resultate aus den Kombinationen \(2-3,\ 1-3,\ 1-2\) seien \(A,B,C\), zugleich die Winkel des durch sie gebildeten Dreiecks, und die ihnen gegenüberliegenden Seiten seien \(a,b,c\).
\begin{center}
\begin{tikzpicture}[x=1cm,y=1cm,line width=.6pt,font=\normalsize]
\coordinate (one) at (0,1.4);\coordinate (two) at (9.4,6.045);\coordinate (three) at (10.5,-.06);
\coordinate (B) at (4,2.8);\coordinate (C) at (7,3.85);\coordinate (A) at (5.25,2.25);\coordinate (P) at (5.44,2.86);
\draw (one)--(B)--(C);\draw (A)--(C)--(two);\draw (B)--(A)--(three);
\draw[line width=.4pt] (5.30155,3.25555)--(P)--(5.6573,2.62236) (P)--(5.1843,2.2789);
\foreach \p in {one,two,three,P}{\fill[white] (\p) circle(.045);\draw (\p) circle(.045);}
\node[above] at (one) {$1$};\node[right] at (two) {$2$};\node[right] at (three) {$3$};
\node[below] at (A) {$A$};\node[above left] at (B) {$B$};\node[right] at (C) {$C$};\node[above right] at (P) {$P$};
\node at (5.48,3.54) {$a$};\node at (6.46,2.92) {$b$};\node at (4.62,2.31) {$c$};
\node at (5.18,3.05) {$x$};\node at (5.49,2.61) {$y$};\node at (5.10,2.65) {$z$};
\end{tikzpicture}
\end{center}
